\section{Theoretical Foundations}
\begin{frame}{Foundations: Performance Difference Lemma}
\begin{itemize}
    \item The identity on the previous slide has a name: the \textbf{Performance Difference Lemma}\footnotemark[1]
    $$\mathcal{J}(\theta) - \mathcal{J}(\theta_\text{old}) = \mathbb{E}_{\tau \sim \pi_\theta}\!\left[\sum_{t \geq 0} \gamma^t A^{\pi_{\theta_\text{old}}}(s_t, a_t)\right]$$
    \item \textbf{Why this matters:} it makes the surrogate $\mathcal{L}(\theta)$ a \textbf{principled approximation} to the true objective $\mathcal{J}(\theta)$ --- not a heuristic
    \item When $\pi_\theta \approx \pi_{\theta_\text{old}}$, the state distributions induced by the two policies are similar, so we can substitute samples from $\pi_{\theta_\text{old}}$ for samples from $\pi_\theta$ with controlled error
    \item As $\pi_\theta$ drifts away from $\pi_{\theta_\text{old}}$, the approximation degrades. This is the theoretical reason we need a trust region --- not just empirical observation that ``big steps break things''
\end{itemize}
\footnotetext[1]{Lemma originally due to Kakade \& Langford (2002); foundation of TRPO in Schulman et al.\ (2015).}
\end{frame}

\begin{frame}{Foundations: Importance Sampling Variance}
\begin{itemize}
    \item Using data from $\pi_{\theta_\text{old}}$ requires importance sampling. The IS-weighted estimator has variance
    $$\mathrm{Var}\!\left[\tfrac{P(x)}{Q(x)}\, f(x)\right] = \mathbb{E}_{x \sim P}\!\left[\tfrac{P(x)}{Q(x)}\, f(x)^2\right] - \mathbb{E}[f]^2$$
    \item The variance \textbf{explodes} when the ratio $\tfrac{P}{Q}$ is large where $f^2$ is large --- a single bad sample dominates the estimate
    \item \textbf{Concretely:} if $\pi_{\theta_\text{old}}(a|s)$ is tiny where $\pi_\theta(a|s)$ is large, the ratio $r_t(\theta)$ blows up. One rare sample under the old policy gets enormous weight and swings the entire gradient estimate
    \pause
    \item For \textbf{trajectory-level} IS: $\frac{p(\tau|\pi_\theta)}{p(\tau|\pi_{\theta_\text{old}})} = \prod_{t=0}^{T} r_t(\theta)$. Small per-step differences \textbf{compound multiplicatively} over $T$ steps --- the variance is hopeless
    \item TRPO and PPO use \textbf{single-step ratios} and constrain them to stay near $1$ --- introduced next\footnotemark[2]
\end{itemize}
\footnotetext[2]{Full derivation in OpenAI's \href{https://spinningup.openai.com/en/latest/algorithms/trpo.html}{Spinning Up: TRPO} and Schulman et al.\ (2015).}
\end{frame}
